\section{무엇을 고쳤나}

\paragraph{주장 1.} 「검정력 있는 변이체」라는 판정은 설정 격자에 「변이 크기 0」인 칸(널칸)을 넣었나에 달려 있다 --- 널칸은 정의상 원본과 같으므로 「어떤 설정에서도 떨어진다」가 원리상 못 켜진다. 앞 노트의 배선 7 개는 널칸을 넣으면 7 개가 「검정력 있음」이고 널칸을 빼면 1 개만 남으며 나머지 6 개는 「구성상 거짓」이다

\paragraph{주장 2.} 그 비대칭은 「진보 서사」를 만든다 --- 세 사이클을 같은 규약(널칸 제외)으로 다시 재면 구성상 거짓이 7 · 6 · 6 로 거의 평평하고, 널칸을 넣으면 5 · 6 · 7 로 방향이 뒤집힌다

\paragraph{주장 3.} 그러나 「널칸이 구성상 있나」 자체는 진짜 자다 --- 판정식을 손으로 뒤집은 변이체에는 「변이 크기 0」이 원리상 없기 때문이다. 그 몫은 세 사이클에서 5/7 · 6/6 · 7/7 로 실제로 올라간다

\paragraph{주장 4.} 산출물 그래프의 고리(SCC)를 상한 없이 면제하면 「어긋남 0」을 살 수 있다 --- 앞 노트는 고리 밖 어긋남 0 과 고리 안 어긋남 8 을 내고 앞 것만 게재했으며, 앞앞 노트에는 같은 자로 7 이 나왔다. 면제를 끄면 두 수는 [8, 8] 로 같다

\paragraph{주장 5.} 탐색 격자는 이 표준오차로 「최적」을 못 정한다 --- 8 칸 격자에서 2·SE 로 안 갈리는 칸이 7 개다. 격자를 넓히는 대신 이분 대비로 다시 세우면 「base 45 대 나머지」는 판정 자에서 0.109682 $\pm$ 0.043888 (t\_clu 2.499153) 로 갈린다

\paragraph{주장 6.} 그리고 그 이분 대비마저 「자의 사실」이다 --- 같은 대비의 t\_clu 가 병기 자에서 -0.096985 과 0.288452 로 문턱을 못 넘는다

\paragraph{주장 7.} 「무엇이 사는가는 자의 사실이다」라는 앞 노트의 헤드라인은 「차」를 재면 서지 않는다 --- 상호작용의 t\_clu 가 판정 자에서 1.467108 · 챔피언 자에서 2.226044 로 문턱의 양쪽에 떨어지지만, 두 자 「사이의 차」 자체는 0.033131 $\pm$ 0.017732 (t\_clu 1.868389) 로 2·SE 를 못 넘는다. 두 수가 문턱의 양쪽에 있다는 것과 두 수가 서로 다르다는 것은 다른 명제다


\section{한계}

이 스텝은 \textbf{컴파일하지 않았고 보내지 않았다}. 그 사실을 원장과 이 파일에 적는다 --- 「안 했다」는 「없다」가 아니다.
